\section{Ablation Study on Local Cache}
\label{app:local}

To validate the necessity of local cache within WG-KV, we conducted an ablation study by removing the local cache. Specifically, we retrained the model using the same objective described in \cref{subsec:objective}, but constrained the local window size $W_{\text{local}}$ to 1. This configuration forces the model to rely exclusively on the learnable gate to decide whether to retain any token \textit{immediately upon generation}, effectively eliminating the grace period for recent tokens.

\cref{fig:threshold_cmp} compares the full WG-KV system against the variant without local cache (``WG-KV w/o Local Cache''). We plot the average distillation loss $\mathcal{L}_{\text{distill}}$ on a held-out validation set of 750 FineWeb-Edu samples against the normalized KV cache size.

The results highlight a sharp degradation in the variant without local cache as admission budget tightens. This empirical evidence strongly substantiates the ``transient utility'' hypothesis discussed in \cref{subsec:sparsity}: recent tokens, even those with negligible long-term utility, attract dense local attention. Without a dedicated local cache to provide a temporary grace period, the admission policy is forced to make premature retention decisions for these high-saliency states. Consequently, the system fails to decouple immediate context needs from long-term retention, confirming that the dual-cache design---combining a transient local cache with a persistent global cache---is essential for efficient and robust long-context inference.